\chapter{分子结构和分子光谱}
\problem{9.1}{$\mathrm{HBr}$ 分子的远红外吸收光谱是一些 $\Delta\widetilde\nu=16.94\,\mathrm{cm^{-1}}$ 等间隔的光谱线。试求 $\mathrm{HBr}$ 分子的转动惯量及原子核间的距离。已知 $\mathrm H$ 和 $\mathrm{Br}$ 的原子量分别为 1.008 和 79.92。}
\answerarea
\problem{9.2}{$\mathrm{HCl}$ 分子有一个近红外光谱带，其相邻的几条谱线的波数是：
\[
2925.78,\quad2906.25,\quad2865.09,\quad2843.56,\quad2821.49\;\mathrm{cm^{-1}}.
\]
$\mathrm H$ 和 $\mathrm{Cl}$ 的原子量分别是 1.008 和 35.46。试求这个谱带的基线波数 $\widetilde\nu_0$ 和这种分子的转动惯量。}
\answerarea\clearpage
\problem{9.3}{$\mathrm{Cl}$ 原子的两同位素 ${}^{35}\mathrm{Cl}$ 和 ${}^{37}\mathrm{Cl}$ 分别与 $\mathrm H$ 原子化合成两种分子 $\mathrm H{}^{35}\mathrm{Cl}$ 和 $\mathrm H{}^{37}\mathrm{Cl}$。试求这两种分子的振动光谱中相应光谱带基线的频率 $\nu_0$ 之比。}
\answerarea
\problem{9.4}{试证明双原子分子各相邻振动能级之间跃迁时发射光的频率与两核间固有振动频率一致，假设两原子间相互作用力为弹性力。}
\answerarea\clearpage
\problem{9.5}{怎样解释分子的组合散射有下列两个特点：
\begin{enumerate}
\item 波长短的伴线比波长长的伴线的强度弱；
\item 随散射体温度的升高，波长短的伴线强度明显增强而波长长的伴线的强度几乎不变。
\end{enumerate}}
\answerarea
\problem{9.6}{光在 $\mathrm{HF}$ 分子上组合散射使某谱线产生波长为 $2670\,\angstrom$ 和 $3430\,\angstrom$ 两条伴线。试由此计算该分子的振动频率和两原子间所作用的准弹性力的弹性系数 $k$ 值。已知 $\mathrm H$ 和 $\mathrm F$ 的原子量分别为 1.008 和 19.00。}
\answerarea\clearpage
